% GERADO por tools/generation/gerar_secoes_latex.py a partir de
% artifacts/paper/secoes/06-ontompo-revisada.md. Nao editar a mao: a proxima geracao
% desfaz. Corrija o Markdown e regere.
O artefato revisado reúne 44 classes, 35 relações e 21 generalizações, com duas partições disjuntas e completas, e acompanha as três dimensões do modelo de referência (Figura 1). Todas as classes chegam definidas, o que o artefato analisado não fazia: 24 definições vêm do glossário do apêndice da conceituação, cinco foram redigidas a partir do papel que o conceito exerce nos diagramas, e quinze nasceram com os conceitos que as rodadas introduziram.

\subsection{As três camadas}

Na camada de \textbf{agentes}, \texttt{Agent} deixou de fornecer princípio de identidade e passou a «category» não-sortal, partida em \texttt{PhysicalAgent} e \texttt{SocialAgent}. Sob a primeira está \texttt{Person}, e sob ela o «role» \texttt{ObservatoryUser}, de quem consulta ou alimenta o observatório; sob a segunda, \texttt{Organization}. \texttt{StakeHolder} é «roleMixin», porque a definição de parte interessada abrange indivíduo, grupo e organização --- três princípios de identidade --- e nenhum sortal único a comporta. \texttt{System}, o ator não-humano que troca dados com o observatório, é papel de \texttt{ComputationalSystem}, o tipo que lhe dá identidade e que a formalização deixara implícito; \texttt{ObservatoryGroup} é «collective».

Na camada de \textbf{estruturas}, \texttt{Software} é «kind» e se especializa em três subtipos mutuamente disjuntos: \texttt{ProjectObservatory}, o observatório propriamente dito, é «subkind»; \texttt{DataManager} responde pela coleta, pelo tratamento, pelo armazenamento e pela disponibilização dos dados; e \texttt{View} é a interface pela qual o observatório se apresenta. Os dois últimos são partes do observatório por «componentOf» e se especializam nos papéis que o modelo distingue: o gerenciamento em \texttt{Collector}, \texttt{Processor}, \texttt{Storer} e \texttt{CrudRepository}; a interface em \texttt{Disseminator}, que divulga dados e análises, \texttt{Reporter}, que coordena a interação entre usuários, projetos e sistemas, e \texttt{CrudView}, que dispara operações sobre o repositório. Os conteúdos são \texttt{Project}, \texttt{Knowledge} --- o que a observação produz --- e \texttt{DataSource}; a infraestrutura reúne \texttt{Hardware}, \texttt{Network} e \texttt{Service}. Sete relatores reificam os vínculos que a prosa enuncia por verbo: \texttt{ProjectDataManagement} liga um gerenciamento ao projeto por cujos dados ele responde e \texttt{ViewProvision}, à interface que ele disponibiliza; \texttt{SocialInteraction} e \texttt{Log} ligam um agente à interface com que ele interage e à que registra sua ação; e \texttt{SoftwareExecution}, \texttt{ServiceProvision} e \texttt{Connection} cobrem execução, serviço e rede.

A camada de \textbf{processos} é a que não existia. \texttt{EtlProcess} é «event» complexo com \texttt{Extract}, \texttt{Transform} e \texttt{Load} como partes, cada uma com participantes declarados: a fonte e o papel de coleta na extração, os de coleta e processamento na transformação, os de processamento e armazenamento na carga. \texttt{CrudOperation} é evento partido em \texttt{CreateOperation}, \texttt{ReadOperation}, \texttt{UpdateOperation} e \texttt{DeleteOperation}, e \texttt{Observation} é o evento pelo qual um agente observa um conteúdo divulgado, do qual nasce o \texttt{Knowledge} correspondente. É essa camada que dá ao artefato um \textit{quando} e um \textit{quem}, e dela dependem quatro das sete questões de competência (\S{}7).

\begin{figure}[H]
\centering
% GERADO por tools/generation/gerar_figuras.py a partir de
% artifacts/ontology/rodada-2/ontompo-rodada-2.ontouml.json.
% Nao editar a mao: a proxima geracao desfaz.
\begin{tikzpicture}[
  every node/.style={font=\fontsize{6}{7}\selectfont},
  classe/.style={draw, rectangle, rounded corners=1pt, inner sep=1pt, minimum width=2.28cm, minimum height=0.52cm, align=center},
  gen/.style={draw, -{Triangle[open, length=3pt, width=3pt]}, thin},
  parte/.style={draw, -{Diamond[length=3pt, width=2.5pt]}, thin},
  med/.style={draw, dashed, dash pattern=on 2pt off 1.5pt, thin},
  part/.style={draw, dotted, thin},
  outra/.style={draw, gray, thin},
]
\fill[black!4, rounded corners=2pt] (-1.26,-2.66) rectangle (13.51,0.56);
\node[anchor=north west, font=\fontsize{5}{6}\selectfont\itshape, text=black!55] at (-1.20,0.54) {Agentes (UFO-C)};
\fill[black!4, rounded corners=2pt] (-1.26,-6.80) rectangle (13.51,-2.82);
\node[anchor=north west, font=\fontsize{5}{6}\selectfont\itshape, text=black!55] at (-1.20,-2.84) {Estruturas (UFO-A)};
\fill[black!4, rounded corners=2pt] (-1.26,-9.42) rectangle (13.51,-6.96);
\node[anchor=north west, font=\fontsize{5}{6}\selectfont\itshape, text=black!55] at (-1.20,-6.98) {Processos (UFO-B)};
\node[classe] (nAgent) at (4.90,0.00) {\fontsize{5}{6}\selectfont«category»\\\fontsize{6}{7}\selectfont\textbf{Agent}};
\node[classe] (nCollector) at (0.00,-4.90) {\fontsize{5}{6}\selectfont«role»\\\fontsize{6}{7}\selectfont\textbf{Collector}};
\node[classe] (nComputationalSystem) at (7.35,-1.52) {\fontsize{5}{6}\selectfont«kind»\\\fontsize{6}{7}\selectfont\textbf{ComputationalSystem}};
\node[classe] (nConnection) at (9.80,-4.14) {\fontsize{5}{6}\selectfont«relator»\\\fontsize{6}{7}\selectfont\textbf{Connection}};
\node[classe] (nCreateOperation) at (7.35,-8.28) {\fontsize{5}{6}\selectfont«event»\\\fontsize{6}{7}\selectfont\textbf{CreateOperation}};
\node[classe] (nCrudOperation) at (7.35,-7.52) {\fontsize{5}{6}\selectfont«event»\\\fontsize{6}{7}\selectfont\textbf{CrudOperation}};
\node[classe] (nCrudRepository) at (0.00,-5.66) {\fontsize{5}{6}\selectfont«role»\\\fontsize{6}{7}\selectfont\textbf{CrudRepository}};
\node[classe] (nCrudView) at (12.25,-4.90) {\fontsize{5}{6}\selectfont«role»\\\fontsize{6}{7}\selectfont\textbf{CrudView}};
\node[classe] (nDataManager) at (2.45,-4.14) {\fontsize{5}{6}\selectfont«role»\\\fontsize{6}{7}\selectfont\textbf{DataManager}};
\node[classe] (nDataSource) at (12.25,-3.38) {\fontsize{5}{6}\selectfont«kind»\\\fontsize{6}{7}\selectfont\textbf{DataSource}};
\node[classe] (nDeleteOperation) at (7.35,-9.04) {\fontsize{5}{6}\selectfont«event»\\\fontsize{6}{7}\selectfont\textbf{DeleteOperation}};
\node[classe] (nDisseminator) at (7.35,-4.90) {\fontsize{5}{6}\selectfont«role»\\\fontsize{6}{7}\selectfont\textbf{Disseminator}};
\node[classe] (nEtlProcess) at (0.00,-7.52) {\fontsize{5}{6}\selectfont«event»\\\fontsize{6}{7}\selectfont\textbf{EtlProcess}};
\node[classe] (nExtract) at (0.00,-8.28) {\fontsize{5}{6}\selectfont«event»\\\fontsize{6}{7}\selectfont\textbf{Extract}};
\node[classe] (nHardware) at (4.90,-3.38) {\fontsize{5}{6}\selectfont«kind»\\\fontsize{6}{7}\selectfont\textbf{Hardware}};
\node[classe] (nKnowledge) at (12.25,-6.42) {\fontsize{5}{6}\selectfont«kind»\\\fontsize{6}{7}\selectfont\textbf{Knowledge}};
\node[classe] (nLoad) at (4.90,-8.28) {\fontsize{5}{6}\selectfont«event»\\\fontsize{6}{7}\selectfont\textbf{Load}};
\node[classe] (nLog) at (9.80,-5.66) {\fontsize{5}{6}\selectfont«relator»\\\fontsize{6}{7}\selectfont\textbf{Log}};
\node[classe] (nNetwork) at (7.35,-3.38) {\fontsize{5}{6}\selectfont«kind»\\\fontsize{6}{7}\selectfont\textbf{Network}};
\node[classe] (nObservation) at (12.25,-7.52) {\fontsize{5}{6}\selectfont«event»\\\fontsize{6}{7}\selectfont\textbf{Observation}};
\node[classe] (nObservatoryGroup) at (12.25,-1.52) {\fontsize{5}{6}\selectfont«collective»\\\fontsize{6}{7}\selectfont\textbf{ObservatoryGroup}};
\node[classe] (nObservatoryUser) at (0.00,-2.28) {\fontsize{5}{6}\selectfont«role»\\\fontsize{6}{7}\selectfont\textbf{ObservatoryUser}};
\node[classe] (nOrganization) at (4.90,-1.52) {\fontsize{5}{6}\selectfont«kind»\\\fontsize{6}{7}\selectfont\textbf{Organization}};
\node[classe] (nPerson) at (0.00,-1.52) {\fontsize{5}{6}\selectfont«kind»\\\fontsize{6}{7}\selectfont\textbf{Person}};
\node[classe] (nPhysicalAgent) at (0.00,-0.76) {\fontsize{5}{6}\selectfont«category»\\\fontsize{6}{7}\selectfont\textbf{PhysicalAgent}};
\node[classe] (nProcessor) at (2.45,-4.90) {\fontsize{5}{6}\selectfont«role»\\\fontsize{6}{7}\selectfont\textbf{Processor}};
\node[classe] (nProject) at (12.25,-5.66) {\fontsize{5}{6}\selectfont«kind»\\\fontsize{6}{7}\selectfont\textbf{Project}};
\node[classe] (nProjectDataManagement) at (2.45,-5.66) {\fontsize{5}{6}\selectfont«relator»\\\fontsize{6}{7}\selectfont\textbf{ProjectDataManagement}};
\node[classe] (nProjectObservatory) at (0.00,-4.14) {\fontsize{5}{6}\selectfont«subkind»\\\fontsize{6}{7}\selectfont\textbf{ProjectObservatory}};
\node[classe] (nReadOperation) at (9.80,-8.28) {\fontsize{5}{6}\selectfont«event»\\\fontsize{6}{7}\selectfont\textbf{ReadOperation}};
\node[classe] (nReporter) at (9.80,-4.90) {\fontsize{5}{6}\selectfont«role»\\\fontsize{6}{7}\selectfont\textbf{Reporter}};
\node[classe] (nService) at (9.80,-3.38) {\fontsize{5}{6}\selectfont«kind»\\\fontsize{6}{7}\selectfont\textbf{Service}};
\node[classe] (nServiceProvision) at (12.25,-4.14) {\fontsize{5}{6}\selectfont«relator»\\\fontsize{6}{7}\selectfont\textbf{ServiceProvision}};
\node[classe] (nSocialAgent) at (4.90,-0.76) {\fontsize{5}{6}\selectfont«category»\\\fontsize{6}{7}\selectfont\textbf{SocialAgent}};
\node[classe] (nSocialInteraction) at (7.35,-5.66) {\fontsize{5}{6}\selectfont«relator»\\\fontsize{6}{7}\selectfont\textbf{SocialInteraction}};
\node[classe] (nSoftware) at (0.00,-3.38) {\fontsize{5}{6}\selectfont«kind»\\\fontsize{6}{7}\selectfont\textbf{Software}};
\node[classe] (nSoftwareExecution) at (7.35,-4.14) {\fontsize{5}{6}\selectfont«relator»\\\fontsize{6}{7}\selectfont\textbf{SoftwareExecution}};
\node[classe] (nStakeHolder) at (9.80,-0.76) {\fontsize{5}{6}\selectfont«roleMixin»\\\fontsize{6}{7}\selectfont\textbf{StakeHolder}};
\node[classe] (nStorer) at (4.90,-4.90) {\fontsize{5}{6}\selectfont«role»\\\fontsize{6}{7}\selectfont\textbf{Storer}};
\node[classe] (nSystem) at (7.35,-2.28) {\fontsize{5}{6}\selectfont«role»\\\fontsize{6}{7}\selectfont\textbf{System}};
\node[classe] (nTransform) at (2.45,-8.28) {\fontsize{5}{6}\selectfont«event»\\\fontsize{6}{7}\selectfont\textbf{Transform}};
\node[classe] (nUpdateOperation) at (12.25,-8.28) {\fontsize{5}{6}\selectfont«event»\\\fontsize{6}{7}\selectfont\textbf{UpdateOperation}};
\node[classe] (nView) at (4.90,-4.14) {\fontsize{5}{6}\selectfont«role»\\\fontsize{6}{7}\selectfont\textbf{View}};
\node[classe] (nViewProvision) at (4.90,-5.66) {\fontsize{5}{6}\selectfont«relator»\\\fontsize{6}{7}\selectfont\textbf{ViewProvision}};
\draw[gen] (nPhysicalAgent) -- (nAgent);
\draw[gen] (nSocialAgent) -- (nAgent);
\draw[gen] (nStakeHolder) -- (nAgent);
\draw[gen] (nPerson) -- (nPhysicalAgent);
\draw[gen] (nOrganization) -- (nSocialAgent);
\draw[gen] (nObservatoryUser) -- (nPerson);
\draw[gen] (nSystem) -- (nComputationalSystem);
\draw[gen] (nCreateOperation) -- (nCrudOperation);
\draw[gen] (nReadOperation) -- (nCrudOperation);
\draw[gen] (nUpdateOperation) -- (nCrudOperation);
\draw[gen] (nDeleteOperation) -- (nCrudOperation);
\draw[gen] (nCrudRepository) -- (nDataManager);
\draw[gen] (nCollector) -- (nDataManager);
\draw[gen] (nProcessor) -- (nDataManager);
\draw[gen] (nStorer) -- (nDataManager);
\draw[gen] (nDisseminator) -- (nView);
\draw[gen] (nReporter) -- (nView);
\draw[gen] (nCrudView) -- (nView);
\draw[gen] (nDataManager) -- (nSoftware);
\draw[gen] (nView) -- (nSoftware);
\draw[gen] (nProjectObservatory) -- (nSoftware);
\draw[parte] (nObservatoryUser) -- (nObservatoryGroup);
\draw[parte] (nStakeHolder) -- (nObservatoryGroup);
\draw[parte] (nSystem) -- (nObservatoryGroup);
\draw[parte] (nDataManager) -- (nProjectObservatory);
\draw[parte] (nView) -- (nProjectObservatory);
\draw[med] (nSocialInteraction) -- (nAgent);
\draw[med] (nSocialInteraction) -- (nReporter);
\draw[med] (nLog) -- (nAgent);
\draw[med] (nLog) -- (nReporter);
\draw[med] (nProjectDataManagement) -- (nDataManager);
\draw[med] (nProjectDataManagement) -- (nProject);
\draw[med] (nViewProvision) -- (nDataManager);
\draw[med] (nViewProvision) -- (nView);
\draw[med] (nSoftwareExecution) -- (nSoftware);
\draw[med] (nSoftwareExecution) -- (nHardware);
\draw[med] (nServiceProvision) -- (nService);
\draw[med] (nServiceProvision) -- (nSoftware);
\draw[med] (nConnection) -- (nHardware);
\draw[med] (nConnection) -- (nNetwork);
\draw[part] (nAgent) -- (nObservation);
\draw[part] (nDisseminator) -- (nObservation);
\draw[part] (nAgent) -- (nCrudOperation);
\draw[part] (nCrudView) -- (nCrudOperation);
\draw[part] (nCollector) -- (nExtract);
\draw[part] (nDataSource) -- (nExtract);
\draw[part] (nCollector) -- (nTransform);
\draw[part] (nProcessor) -- (nTransform);
\draw[part] (nProcessor) -- (nLoad);
\draw[part] (nStorer) -- (nLoad);
\draw[parte] (nExtract) -- (nEtlProcess);
\draw[parte] (nTransform) -- (nEtlProcess);
\draw[parte] (nLoad) -- (nEtlProcess);
\draw[outra] (nKnowledge) -- (nObservation);
\draw[outra] (nSoftware) -- (nHardware);
\draw[gen] (-1.14,-9.80) -- (-0.74,-9.80);
\node[anchor=west, font=\fontsize{5}{6}\selectfont] at (-0.69,-9.80) {generalização};
\draw[parte] (3.70,-9.80) -- (4.10,-9.80);
\node[anchor=west, font=\fontsize{5}{6}\selectfont] at (4.15,-9.80) {«componentOf», «memberOf», «participational»};
\draw[med] (8.55,-9.80) -- (8.95,-9.80);
\node[anchor=west, font=\fontsize{5}{6}\selectfont] at (9.00,-9.80) {«mediation»};
\draw[part] (-1.14,-10.12) -- (-0.74,-10.12);
\node[anchor=west, font=\fontsize{5}{6}\selectfont] at (-0.69,-10.12) {«participation»};
\draw[outra] (3.70,-10.12) -- (4.10,-10.12);
\node[anchor=west, font=\fontsize{5}{6}\selectfont] at (4.15,-10.12) {«material», «creation», «derivation»};
\end{tikzpicture}

\caption{A OntoMPO revisada: as três camadas do modelo, com as classes que cada
uma reúne, suas generalizações e as relações entre elas. Gerada a partir de
\texttt{ontompo-rodada-2.ontouml.json}.}
\label{fig:integrado}
\end{figure}

\subsection{Do modelo em OntoUML à ontologia em OWL}

A implementação não foi traduzida à mão: o modelo revisado passa pela transformação gUFO oficial \cite{almeida2019gufo}, executada pela \texttt{ontouml-js}, e cada elemento da OWL corresponde ao elemento do modelo que o gerou --- é essa correspondência a garantia de preservação semântica que uma tradução manual não oferece.

Sobre o gerado incidem cinco customizações, todas \textbf{aditivas por regra}: nenhuma tripla produzida pela transformação é removida ou alterada, e a contenção do gerado dentro do customizado é conferida a cada execução. A \textbf{C1} declara localmente os termos da gUFO que o artefato referencia --- tipo, rótulo, domínio e alcance ---, para que o arquivo se sustente quando a importação não é resolvida; a \textbf{C2} devolve a cada classe sua definição e seus rótulos preferidos; a \textbf{C3} declara disjuntos os «kind» do modelo, consequência da semântica do estereótipo que a transformação não carrega, e os três subtipos de \texttt{Software}; a \textbf{C4} dá nome à inversa de cada propriedade de objeto, sem o que a proveniência teria de ser percorrida por caminho invertido a cada passo; e a \textbf{C5} acrescenta título, licença e versão. Somam 648 triplas, publicadas à parte como o acréscimo exato.

Três decisões foram deliberadamente \textbf{não} tomadas: os identificadores não foram renomeados, porque isso quebraria a correspondência com o modelo; os papéis irmãos não foram declarados disjuntos, porque nada na fonte impede que um componente colete e armazene; e nenhuma axiomatização pesada foi acrescentada, por decisão de escopo (\S{}8).

\subsection{Métricas do artefato}

A Tabela 2 conta o que está nos arquivos, sem a gUFO importada.

\begin{table}[H]
\centering
\caption{Métricas dos artefatos OWL.}
\label{tab:2}
\scriptsize
\begin{tabular}{p{\dimexpr0.3421\linewidth-2\tabcolsep\relax}>{\centering\arraybackslash}p{\dimexpr0.1711\linewidth-2\tabcolsep\relax}>{\centering\arraybackslash}p{\dimexpr0.2105\linewidth-2\tabcolsep\relax}>{\centering\arraybackslash}p{\dimexpr0.2763\linewidth-2\tabcolsep\relax}}
\hline
 & \textbf{linha de base} & \textbf{revisada, gerada} & \textbf{revisada, customizada} \\
\hline
classes nomeadas & 32 & 44 & 44 \\
propriedades de objeto & 28 & 34 & 68 \\
propriedades de dados & 0 & 0 & 0 \\
subsunções & 60 & 86 & 86 \\
conjuntos de disjunção & 0 & 2 & 4 \\
pares de inversas nomeadas & 0 & 0 & 34 \\
axiomas & 381 & 558 & 809 \\
anotações & 60 & 78 & 475 \\
triplas & 441 & 636 & 1284 \\
\hline
\end{tabular}
\end{table}

Duas linhas pedem leitura. As \textbf{propriedades de dados são zero} nas três colunas, e é número para declarar, não para esconder: a conceituação especifica atributos para cada conceito, e nenhum chegou ao diagrama (\S{}8). O salto nas \textbf{anotações} é a C2 medida.
